% One-page resume: Aakarsh Anand. The full CV (everything, no page limit) is cv.tex.
% Build: make -C cv

\documentclass[10pt,letterpaper]{article}
\usepackage[margin=0.5in,top=0.45in,bottom=0.4in]{geometry}
% Shared look and header for resume.tex and cv.tex. Change styling or contact details here only.
% EB Garamond for reading text; Gill Sans (a macOS system font) for the name, headings, dates, and metadata.
% No italics anywhere; venues are bold.

\usepackage[lining,scaled=1.06]{ebgaramond}
\setsansfont{Gill Sans}
\usepackage{microtype}
\usepackage{xcolor}
\usepackage{enumitem}
\usepackage{titlesec}
\usepackage{textcase}
\usepackage{hyperref}

\definecolor{accent}{HTML}{3A6532}
\definecolor{muted}{HTML}{5F665F}
\definecolor{rulecolor}{HTML}{BCCFB5}

\hypersetup{colorlinks=true, urlcolor=accent, linkcolor=accent, pdfauthor={Aakarsh Anand}}

\pagestyle{empty}
\hyphenation{RECOMB NeurIPS}
\setlength{\parindent}{0pt}
\setlength{\parskip}{0pt}

% Section headings: small letterspaced sans capitals over a hairline
\newcommand{\secnote}{}
\newcommand{\sectionnote}[1]{\gdef\secnote{#1}}
\newcommand{\headwithrule}[1]{{\addfontfeature{LetterSpace=12}\MakeTextUppercase{#1}}%
  \hspace{0.7em}{\color{rulecolor}\leaders\hrule height 0.72ex depth -0.6ex\hfill}%
  \ifx\secnote\empty\else\hspace{0.7em}{\normalfont\sffamily\footnotesize\color{muted}\secnote}\global\let\secnote\empty\fi}
\titleformat{\section}[block]{\sffamily\color{accent}}{}{0pt}{\headwithrule}
\titlespacing*{\section}{0pt}{8pt}{3.5pt}

\setlist[itemize]{leftmargin=1em, labelsep=0.5em, itemsep=1.5pt, parsep=0pt, topsep=2pt, partopsep=0pt,
  label={\raisebox{0.15ex}{\scriptsize\color{muted}\textbullet}}}

% ---------- macros ----------
\newcommand{\me}{\textcolor{black}{\textbf{A.~Anand}}}
\newcommand{\dates}[1]{{\sffamily\small\color{muted}#1}}
\newcommand{\note}[1]{{\color{muted}#1}}
\newcommand{\vn}[1]{{\color{accent}\textbf{#1}}}
% Separator dot: binds to the preceding word, so a line can end with it but never start with it
\newcommand{\sep}{\unskip\nobreak\enspace{\color{muted}\textperiodcentered}\hspace{0.5em}\ignorespaces}
% \role{title}{organization}{dates}
\newcommand{\role}[3]{\textbf{#1}, #2\hfill\dates{#3}\par}
% \thread{name}{optional note}
\newcommand{\thread}[2]{\vspace{3pt}{\sffamily\color{accent}#1}%
  \if\relax\detokenize{#2}\relax\else{\sffamily\small\color{muted}\sep #2}\fi\par}
% \pub{title}{authors}{venue}{url}: title on one line, authors and venue beneath
\newcommand{\pub}[4]{\item \href{#4}{\textcolor{black}{#1}}\\{\small\color{muted}#2}\sep{\small #3}}

% ---------- header ----------
\newcommand{\header}{%
  \begin{minipage}[b]{0.5\textwidth}
    {\sffamily\fontsize{23}{25}\selectfont\color{accent}\addfontfeature{LetterSpace=9}AAKARSH ANAND}
  \end{minipage}%
  \begin{minipage}[b]{0.5\textwidth}
    \raggedleft\sffamily\small
    \href{mailto:aakarsh@cs.ucla.edu}{aakarsh@cs.ucla.edu}\sep
    \href{https://aakarsh-anand.github.io}{aakarsh-anand.github.io}\\[1pt]
    \href{https://scholar.google.com/citations?user=icSvTl8AAAAJ}{Google Scholar}\sep
    \href{https://github.com/aakarsh-anand}{GitHub}\sep
    \href{https://www.linkedin.com/in/aakarsh-anand-ml/}{LinkedIn}
  \end{minipage}\par
  \vspace{4pt}
  {\color{muted}Machine learning researcher working on foundation models and statistical genetics for
  health data}\par
  \vspace{-2pt}}

\hypersetup{pdftitle={Aakarsh Anand - Resume}}

\begin{document}

\header

% ---------- education ----------
\section{Education}
\role{Ph.D., Computer Science}{UCLA \sep Major field: AI \sep Advisor: Sriram Sankararaman}{2024 -- 2029 (expected)}
\role{M.S., Computer Science}{UCLA \note{(earned during the Ph.D.)}}{2026}
\role{B.S., Computer Science}{UCLA}{2020 -- 2024}

% ---------- research ----------
\section{Research Experience}
\role{Graduate Student Researcher}{Sankararaman Lab, UCLA \note{(undergraduate researcher 2022--24)}}{Jan 2022 -- present}

\thread{Foundation models for wearable sensors}{with Yuzhe Yang's lab}
\begin{itemize}
  \item Co-first author of \textbf{Inertia-1} (NeurIPS 2026). Pretrained motion foundation models (5M--100M params,
    6 architectures) on 18.2M~hours of accelerometry from 115K+~people and compared data, model,
    and training choices on 15 datasets.
  \item Designed and ran the sensor-fusion study, which found that multi-sensor models consistently beat single-sensor ones.
  \item Co-first author of the follow-up journal paper (in prep.)\ on diagnosis, prognosis, and disease progression in
    five cohorts, including the genetics of the learned embeddings and early work pairing them with LLMs.
\end{itemize}

\thread{Deep learning for genetics}{}
\begin{itemize}
  \item Developing a training loss that makes a network's learned trait heritable and genetically correlated with
    disease, so a GWAS on it can find loci the original measurement misses. It works with any architecture.
\end{itemize}

\thread{Statistical genetics at biobank scale}{}
\begin{itemize}
  \item Leading a genome-wide test for gene--gene interactions in rare disease, supported by an Anthropic grant.
  \item Leading a single-cell model of how genetic effects on gene expression shift across cell types and states.
  \item Co-leading an atlas of genetic correlation across 3,570 UK Biobank trait pairs and 53 tissues.
  \item Ran simulations (calibration, power, runtime) and real-data analyses for the papers below.
\end{itemize}

% ---------- publications ----------
\sectionnote{* equal contribution}\section{Publications}
\begin{itemize}
  \pub{Inertia-1: An Open Exploration of Wearable Motion Foundation Models}
      {Z.~Xu*, \me*, S.~Jiang*, C.~Zhuang, Z.~Shuai, S.~Sankararaman, Y.~Yang}
      {\vn{NeurIPS 2026}\sep\href{https://github.com/yang-ai-lab/Inertia-1}{code}}
      {https://arxiv.org/abs/2607.06617}
  \pub{A biobank-scale method for learning modulators of gene-environment interaction underlying human complex
      traits from multiple environmental exposures}
      {Z.~Liu, A.~Ramteke, \me, A.~Gorla, M.~Jeong, S.~Sankararaman}
      {\vn{Genome Research 2026}\sep selected from RECOMB 2026}
      {https://doi.org/10.1101/gr.282105.126}
  \pub{A biobank-scale test of marginal epistasis reveals genome-wide signals of polygenic interaction effects}
      {B.~Fu, A.~Pazokitoroudi, Z.~Shi, A.~Kar, A.~Xue, \me, P.~Anand, et al.}
      {\vn{Nature Genetics 2025}}
      {https://doi.org/10.1038/s41588-025-02411-y}
  \pub{A scalable adaptive quadratic kernel method for interpretable epistasis analysis in complex traits}
      {B.~Fu*, P.~Anand*, \me*, J.~Mefford, S.~Sankararaman}
      {\vn{Genome Research 2024}\sep selected from RECOMB 2024}
      {https://doi.org/10.1101/gr.279140.124}
  \pubnolink{Disentangling phenotype scale, amplification, and effect heterogeneity underlying polygenic GxE
      architecture of complex traits}
      {Z.~Liu, \me, X.~Qie, M.~Costantino, M.~Jeong, A.~Gorla, N.~Zaitlen, A.~Dahl, S.~Sankararaman}
      {submitted to \vn{Nature Genetics}}
  \pub{Comprehensive gene heritability estimation reveals the genetic architecture of rare coding variants
      underlying complex traits}
      {Z.~Liu, B.~Fu, M.~Jeong, P.~Anand, \me, et al.}
      {bioRxiv 2025\sep in revision at \vn{Nature Genetics}}
      {https://doi.org/10.1101/2025.10.07.681018}
\end{itemize}

% ---------- awards ----------
\section{Awards and Funding}
\role{Anthropic AI for Science Grant}{Rare Disease Research}{2026}
\role{NIH T32 Predoctoral Fellowship}{UCLA Genomic Analysis Training Program}{2026 -- 2027}
\role{Warren Alpert Computational Biology and AI Network Fellowship}{UCLA}{2024}

% ---------- teaching ----------
\section{Teaching}
\role{Teaching Assistant}{UCLA}{2025 -- 2026}
\begin{itemize}
  \item AI for Computational Genomics (Spring 2025, $\sim$120 students). Co-wrote three new assignments
    for its AI redesign.
  \item Machine Learning Applications in Genetics (Spring 2026) and Machine Learning Applications
    in Biomedicine (Winter 2025).
\end{itemize}

% ---------- earlier experience ----------
\section{Earlier Experience}
\role{Software Engineering Intern}{LabCorp \note{\sep Snakemake pipelines for high-volume sequencing data}}{Jun -- Sep 2023}
\role{Research Assistant}{Boutros Lab, UCLA \note{\sep cancer genomics pipelines, co-author of
  \href{https://doi.org/10.1016/j.crmeth.2026.101340}{Metapipeline-DNA}}}{2021 -- 2022}

% ---------- skills ----------
\section{Skills}
\textbf{Machine learning:} PyTorch, scikit-learn, Transformers, self-supervised pretraining, time-series and wearable sensor data\par
\textbf{Statistics and genetics:} variance-component models, heritability and genetic correlation, GWAS, LDSC, single-cell genomics\par
\textbf{Engineering:} Python, NumPy/pandas, Weights \& Biases, SLURM, Claude Code \sep
\textbf{Data:} UK Biobank, NHANES, OneK1K, GTEx\par

\end{document}
